\documentclass[a4paper,11pt,openany]{ltjsbook}
\usepackage{luatexja}
\usepackage{luatexja-fontspec}
\usepackage{amsmath,amssymb,amsfonts,amsthm}
\usepackage{geometry}
\geometry{top=24mm,bottom=24mm,left=24mm,right=24mm,footskip=10mm}
\usepackage{tcolorbox}
\tcbuselibrary{skins,breakable}
\usepackage{hyperref}
\usepackage{fancyhdr}
\usepackage{graphicx}
\usepackage{tikz}
\usetikzlibrary{arrows.meta,calc,positioning,shapes.geometric}

\definecolor{primaryblue}{RGB}{20, 60, 110}
\definecolor{darkslate}{RGB}{35, 45, 60}
\definecolor{accentcrimson}{RGB}{160, 30, 45}
\definecolor{softblue}{RGB}{238, 244, 252}

\hypersetup{
    colorlinks=true,
    linkcolor=primaryblue,
    citecolor=darkslate,
    urlcolor=primaryblue,
    pdftitle={アルキメデス —— 古代幾何学・力学の巨人と近代科学の源流},
    pdfauthor={サー・トーマス・リトル・ヒース / 教育情報分析研究会}
}

\pagestyle{fancy}
\fancyhf{}
\fancyhead[LE,RO]{\small\thepage}
\fancyhead[LO]{\small\rightmark}
\fancyhead[RE]{\small\leftmark}
\renewcommand{\headrulewidth}{0.4pt}
\setlength{\headheight}{18pt}

\begin{document}

\frontmatter

\begin{titlepage}
    \centering
    \vspace*{2.0cm}
    {\Huge\textbf{アルキメデス}\par}
    \vspace{0.4cm}
    {\Large\textbf{古代幾何学・力学の巨人と近代科学の源流}\par}
    \vspace{0.8cm}
    {\large — — 原著完全翻訳・数学校閲決定版 — —\par}
    \vspace{1.5cm}
    {\large 原題：\textit{Archimedes} (Pioneers of Progress: Men of Science, 1920)\par}
    \vspace{0.3cm}
    {\large Project Gutenberg eBook No. 35550\par}
    \vspace{2.5cm}
    {\Large サー・トーマス・リトル・ヒース（Sir Thomas Little Heath）著\par}
    \vspace{0.5cm}
    {\large 教育情報分析研究会 訳・数理歴史校閲\par}
    \vfill
    {\large 2026年10月 完全決定版\par}
    \vspace*{1.0cm}
\end{titlepage}

\tableofcontents
\clearpage

\chapter*{解題：『アルキメデス』の学術的意義と現代的価値}
\addcontentsline{toc}{chapter}{解題：『アルキメデス』の学術的意義と現代的価値}

本書は、古代ギリシャ数学・天文学史の最高権威であるサー・トーマス・リトル・ヒース（Sir Thomas Little Heath, 1861--1940）が、ロンドンのキリスト教知識普及協会（SPCK）の叢書「進歩の先駆者たち：科学の巨人（Pioneers of Progress: Men of Science）」の一冊として1920年に著した『\textit{Archimedes}』の完全日本語翻訳・数学校閲版である。

ヒースはケンブリッジ大学トリニティ・カレッジで古典学と数学を修め、英国財務省の高級官僚（常任次官）を務める傍ら、ユークリッド『原論』全13巻の英訳注解や『ギリシャ数学史（A History of Greek Mathematics）』など、20世紀の西洋古典科学史を決定づける不朽の記念碑的業績を遺した。その浩瀚な学識をコンパクトな一冊に結晶化させたのが本書である。

アルキメデス（Archimedes of Syracuse, ca. 前287--前212）は、単なる「てこの原理」や「エウレカ（王冠の比重測定）」の伝説的発明家にとどまらない。彼はユークリッドやエウドクソスの幾何学を継承しつつ、現代の積分学の直観的先駆けとなった「力学的求積法（Method of Mechanical Theorems）」を考案し、厳密な「取り尽くし法（Method of Exhaustion）」によって球・円柱・回転放物面・らせん・放物線の面積や体積を正確に計算し尽くした、人類史上最も傑出した数学者の一人である。

本書は、シラクサ防衛戦や彼の劇的な生涯（第I章）、エウドクソスやヒッポクラテスからアルキメデスに至るギリシャ幾何学の雄大な発展史（第II章）、奇跡的に現存した写本伝承（第III章）、円周率の厳密な上下不等式評価や球と円柱の驚異の幾何学（第IV章）、宇宙を埋め尽くす砂の粒の数を指数的命数法で数え上げた『砂粒論』（第V章）、そして重心論・てこの原理（第VI章）や流体静力学・浮体の安定条件（第VII章）に至るまで、アルキメデスの知的宇宙の全貌を余すところなく明解に解き明かしている。

高校・大学初年級の数学・物理知識（三角法、初等幾何、力学の基礎）があれば、2000年以上前の天才がいかにして現代の微積分学や理論物理学の礎を独力で築き上げたのか、その息を呑む論理展開を鮮明に追体験することができるであろう。

\clearpage
\mainmatter
\chapter{第1章：アルキメデスの生涯と逸話}
\section*{ARCHIMEDES}

第一章。

アルキメデス。

一般の人々がアルキメデスについて即答で述べる内容を尋ねられれば、おそらく最も多くは彼に関する有名な逸話の一つか二つを引用できるだけである：浴槽にいるときに何かの問題の解決策を見つけた後、その喜びに満ちて裸で家へ走り、[ギリシャ語: エウレカ、エウレカ]（または我々は「それがわかった、それがわかった」と言う）と叫んだこと；あるいは「私が立つ場所を教えてくれれば地球も動かすことができる」などと言ったこと；また第二次パルティック戦争でシラクサが包囲され、その城塞が陥落した際に、地に描いた図案から立ち去るよう命じられたローマ兵によって殺害されたこと。

歴史と伝承では、アルキメデスは古代最大の数学者であり、おそらく人類で最も優れた数学的天才であるにもかかわらず、その原始的な数学発見に関する詳細については、数学の専門家以外にはほとんど知られていないことが懸念される。

アルキメデスは主に多くの巧妙な機械装置の発明者として歴史と伝統で知られている。これらの発明は自然に一般の人々の想像力を引きつけるものであり、純粋数学の奥深さよりも魅力的である。

アルキメデスに関するほとんど全ての情報はマルケルスによるシラクサ包囲戦について記したポリビウスとプラUTERSの記述を通じて伝わっている。彼は紀元前212年にシラクサが陥落した際に没しており、当時75歳前後であったことから、彼は紀元前287年頃に生まれたと考えられる。彼は天文学者フィディアスの息子であり、シラクサの王ヒロオンとその息子ゲロンと親友兼姻族関係にあった。彼はアレクサンドリアで Euclid の後継者たちのもとで学び、紀元前300年頃 Euclid がまだ生きていたとは限らない。おそらくアレクサンドリアではコンオン・オブ・サモスとの出会いがあり、彼はその才能を称賛していた。

アルキメデスは、エラトステネスと共に友人として敬われた数学者であり、彼はシラクサに戻ってから数学的研究に完全に没頭する生活を送った。偶然にも、彼の巧妙な機械的発明により有名になったが、これらは彼にとって「幾何学の遊び」の一形態で、実用的な価値には重きを置かなかった。

プラUTERSの言葉によれば、「彼は高い精神を持ち、深い魂を持っており、科学知識に富んでいたため、これらの発明が人間を超えた洞察力をもたらしたとしても、書籍でそのような分野について記述することには同意しなかった。彼は機械工学や実用的な技術を軽蔑し、これらは日常の需要と混ざる美しさと精巧さの探求にのみ野心を置いた」。

シラクサの包囲戦中、アルキメデスはローマの包囲軍に対してあらゆる種類の機械を考案した。投石機は巧妙に設計されており、遠距離でも近距離でも使用可能であり、壁に開けられた穴から矢弾を放つ装置もあった。他の装置には、壁外側に延びる長い動く棒があり、これらの棒は敵の船や彼らが「サムブカ」と呼ぶ、その名前の似た保護された梯子に重い物体を落とした。

このようにして、アルキメデスは数学的研究と包囲戦での機械的発明の両方について記述することには同意しなかった。彼はこれらの実用的な技術を軽蔑し、美しさと精巧さに焦点を当てた探求のみに野心を置いたのである。

アルキメデスは数学的な考察に没頭して生きて、そして死んだ。彼の死に関する記述にはいくつかの詳細が異なっている。プラUTERSは次のように述べている。「マルケルスはアルキメデスの死を最も悲しんだ。実際、彼は何かの図形問題を考え込んでおり、その研究に没頭して目と心を一心に注いでいたため、ローマ軍の侵入や城郭の包囲などを見逃した。そして、兵士が突然近づき、マルケルスに従うよう命じたとき、彼は問題を証明するまでその研究から離れないと言った。これに対して兵士は激怒し、剣を抜いて彼を殺した。またの説では、ローマ人は

アルキメデスの生涯と逸話には、彼が円筒の中に球を収める図形を描くように命じられた兵士に「友よ、図形から離れろ」と言ったという最も風景的な説明がある。その兵士は激怒して彼を殺した。

アルキメデスは、円筒と内接する球の比率（3/2）が自身の最大の成績であると自認し、友人や親族にその墓碑に円筒と球の図形を刻むよう求めたという。シケル島で執事官だったキセロンは、この墓標を修復した。

アルキメデスの生涯に関する以上の特定情報以外には、彼の人格について我々が好まずに変更するであろうと思われる数々の逸話しかない。彼が抽象的な学問に完全に没頭していたことを示す逸話では、彼は食事や生活必需品を忘れることがあり、火の灰に幾何学的図形を描いたり、自身を油で身を包んだときにその上に図形を描くこともあったという。

アルキメデスは、ヒロンから金の冠に銀が含まれていないかを調べるよう命じられた問題の解決策を見つけるために浴槽に入っていたときにその解答を見つけ出し、その結果裸で町を走り、家へ向かい、「エウレカ、エウレカ」と叫んだとされる。

アルキメデスは、与えられた重さを与えられた力で動かす問題の解決策を見つけることに結びつけた言葉として有名な「動かすための場所を与えれば、地球自体も動かせる」([ギリシャ語: ドス・マイ・ポウ・スト・カイ・キノ・テン・ゲン]または、一部の報告では、広いドリオ語で[ギリシャ語: パ・ボ・カイ・キノ・タ・ガΝ])を発した。プラUTERSは、アルキメデスがヒロンに対して、与えられた力で何でも動かすことができると宣言し、その論証の妥当性に依拠して、もう一つの地球へ渡りこの地球を動かすことを自慢していたと記述している。また、ある報告では、アルキメデスは王の兵器庫から三重帆船を選んで、多くの人で引き上げたという。その船には多くの乗客と満載の荷物が積まれており、彼自身は遠くに座りながら、大きな力を使わずに静かに複合式の滑車を動かし、船を平穏に安全に引き寄せたとされる。ヒロンは驚き、その問題を実践的に示して見せるようアルキメデスに求めたが、彼は複合式の滑車を使って船を動かしたという。

他には、別の報告では、機械が螺旋([ギリシャ語: ヘリックス])として記述されている。この用語は、円筒形の螺旋でハンドルによって回転させられ、斜めの歯を持つ歯車と噛み合って働く螺線を指すものと考えられる。

アルキメデスは太陽、月、そして五つの惑星の天体運動を模倣する球体の装置も発明した。シーザロはこの装置を見たことがあるとし、その説明の中で、月の周期と太陽の視覚的な運動を非常に正確に表現しており、短い期間では日食や月食まで示すことができると述べている。水力で動かされていた可能性もあるが、パッポスは「球体の製作法を理解し、円形運動を利用した天体模型を作成する人々」について言及している。いずれにせよ、アルキメデスは天文学に関心を持っていたことは確かである。リビウスは彼を「天空と星の唯一の観察者」と呼んでいる。ヒッパルコスによれば、「これらの観測から明らかに年長との差異は非常に小さいが、至点に関してはアルキメデスも私自身も一日の四分の一程度誤っていると考えている。」したがって、アルキメデスは一年の長さについて考慮していたことがわかる。マクロビウスによれば、彼は惑星の距離を発見したとされている。アルキメデス自身は『塵算』において、太陽の視覚的な直径を測定するための装置について述べており、これは眼から見た太陽が生成する角を指すものである。

ローマ船隊を放火するために Burning-glasses や凹面鏡を使用したという話は、シルニウス（2世紀後半）以降の証拠しか見つからない。しかし、反射鏡の形態、例えば回転パラボロイドが存在し、その凹面表面に垂直方向に射入する全ての光線を焦点に集束させる可能性は否定できない。

\clearpage

\chapter{第2章：アルキメデスに至る古代ギリシャ幾何学の歩み}
\section*{GREEK GEOMETRY TO ARCHIMEDES}

第2章。

古代ギリシャ幾何学の歩みとアルキメデスへの至り。

アーキメデスの位置や彼の業績の意義を正しく理解するためには、ギリシャ幾何学がエウクレidesまで遡る発展過程を概観することが必要である。

ギリシャの著者はヘロドトスから下で一致して言っているところによれば、幾何学はエジプト人によって発明され、ギリシャにエジプト経由で伝わったという。ある説では、次のように述べられている。

「多くの人々が幾何学はエジプトで発明されたと主張し、その起源は土地測量のために生じた測量からである。これは毎年夏至に洪水をもたらすナイルの氾濫により、個々の所有者の境界線が消え失ったため必要となった。このような機会から幾何学や他の科学の発見が生まれることは驚くべきことではない。なぜなら、何かが発展的な動きで進化するときには、不完全な状態から完全なものへと自然に移行することがある。感覚から論理へ、そして理解へと移行するのは当然である。フェニキア人においては、商取引を通じて正確な数の知識が生まれたように、エジプト人においても上記の理由により幾何学が発明された。

「テラスは最初にエジプトを訪問し、その後ギリシャへとこの学問を導入した」という説もある。」

しかし、エジプトの幾何学はほとんど実用的であり、土地測量士、農民、商人の要請を超えて発展しなかったことが明確である。彼らは2000BC頃から、3:4:5という辺の比を持つ三角形において、二つの短い辺で囲まれる角が直角であることを見抜いており、実用的な手段としてそのような三角形を使って直角を描画していた。また、円周率や他の測定に関する公式も存在したが、それらは比較的不正確なものであった。

古代ギリシャ幾何学の歩みにおいて、アルキメデスに至るまで、三角形、平行台形、円の面積を算出する方法が知られていた。また、ピラミッドの構築においては、相似直角三角形の概念を使用していた；さらに、「se-qet」という名前で半基底辺の半分に対する高さの比、すなわち斜面角度の余接を表すものがあった。しかし、幾何学における単一の一般定理はエジプト人には見つからなかった。3, 4, 5の三角形が直角であるという知識は、ピタゴラスの有名な定理（ユークリッドI，47）を意味する一般的な命題を示すものではないほど遠くに留まっていた。実際、幾何学の科学自体はまだ発見されていなかった；これは純粋な推測の天才が必要とし、ギリシャ人は世界中のどの民族よりもその才能を持っていた。

テラスはエジプトを訪問し、そこで祭司から教わったことを学んだ後、幾何学をギリシャに導入した。ギリシャの科学と哲学全体がテラスから始まったと言っても過言ではない。彼の生年は紀元前624年頃から547年頃までである。イオニアの哲学者最初の一人であり、582年から581年に第七賢人として宣言された彼は、天文学者、数学者、技術者、政治家、商人としても輝いていた。天文学では、5月28日に太陽食が予測され、四つの天文季節の不均衡を発見し、小熊座を使用して極星を見つける方法として大熊座を使う必要があることを助言した。幾何学においては、彼に帰属される以下の定理があり、その性質はギリシャ人が理論から始めるべき最初のところから始めなければならなかったことを示している－（1）円は直径で二等分される（ユークリッドI，定義17）、（2）等脚三角形の底角は等しい（ユークリッドI，5）、（3）二直線が交わるとき、

古代ギリシャ幾何学の歩みにおいて、アルキメデスに至るまで、垂直対頂角は等しい（ユークリッドI，15），二つの三角形がそれぞれ二つの角と一方の辺を等しくするとき、その三角形はすべてにおいて等しい（ユークリッドI，26）という規則がある。彼は最初に半円の中に入れる直角三角形を描いたとされる：これはつまり、半円の角度が直角であることを初めて発見したということである。また、実用幾何学における二つの問題を解決した：（1）海で船から陸への距離を測る方法を示した（これには上記の（4）の命題が使われたとされる），（2）ピラミッドの高さを地上に投げられた影を利用して測った（これは相似な三角形を利用することを含む）。この手法はエジプト人がピラミッド建設のために使用した方法と同じである。

テレスコルスから、ピタゴラス派が登場する。ピタゴラス派は「数学」という用語（教える対象）を専門的な意味で最初に使用したとされる。彼らはまた、自己のための学問としての数学を追求し、自由教育の一環として数学を導入した。ピタゴラスはメネサルコスの息子であり、紀元前572年頃サモスで生まれ、メタポントムで大いなる年齢（75歳か80歳）で亡くなった。彼の興味はテレスコルスと同様に多岐にわたっており、知識を追求するために行われた旅路はおそらくさらに広範囲であった。テレスコルスと似て、あるいは彼の提案により、エジプトへ訪問し、長期間（22年、一部では言っている）学んだ。

ピタゴラス派の教えから、ピタゴラス自身に帰属すべき部分を剥離するのは困難である。学校のメンバーはすべてを師に帰すという習慣があったためである（[ギリシャ語: autos epha]、ipse dixit）。天文において少なくとも二つは彼に帰属できる：地球は以下のように考えられた。

球体の形状を有し、太陽、月、惑星は日々の回転とは異なる方向に独立した運動を持っていると彼は認識しました。しかし、彼は地球中心説を堅持し、その理論は彼の後継者によって地球が中心に留まらず、他の惑星や太陽、月とともに「中央の火」を中心に回るという考えへと発展しました。

彼の最も驚くべき発見の一つは、音階間隔の依存性が振動弦の長さに由来することでした。オクターブの比は2:1、五度の比は3:2、四度の比は4:3です。算術において彼は、共通比を用いた平均と比例の理論を初めて説明しました。彼は数そのものの性質を考えることで数論の基礎を築き、素数、奇数、偶数などについて考察しました。図形的な数（点で表された四角形、長方形、三角形など）を通じて、数と幾何学との関連性を示しました。これらの数の性質から考えると、ピタゴラス派が「すべてのものを数に例える」ことや、「数の原理はすべてのもののもととなる原理である」という考え方に至ったのは容易に理解できます。

次にピタゴラスの幾何学における業績について述べます。彼がエジプトから帰国し、サモスで学校を設立しようとしたとき、サモアンの人々は冷淡であり、そのため彼自身の幾何学が消滅しないように特別な措置を講じました。ギムナズイオンに赴き、適切な青年を見つけ出し、貧しい彼を雇い、定理を一つ理解するごとに六ペンスを約束することで幾何学を学ばせました。

古代ギリシャ幾何学の歩みは、若者の心を惹きつけ、その本質自体に魅力を感じさせた。ピタゴラスは正しく、六ペンスを手放す代わりに数学を続けることに喜びを見出していたと判断した。彼は自分自身が貧しく生計を得るための努力を余儀なくされていると言明し、その結果、若者は学問を放棄することなく、ピタゴラスから各命題につき六ペンスを支払うことを申し出た。

幾何学において、ピタゴラスは主題の基礎を築くことに自らを捧げ、重要な定義から始め、基本的な原理を探求した。彼に帰属される最も有名な命題は当然、直角三角形において斜辺の平方が他の二辺の平方の和と等しいという定理である（ユークリッドI, 47）。そして、ギリシャ伝統がこの定理の証明を彼に帰属する普遍的な信用を持っていることを考慮に入れると、その伝承は正しいと考えるのが適切だと考えられる。これはまた別の伝承によって部分的に確認され、それはピタゴラスが二つの数を見つける一般式を見つけたことである。その数の平方和が平方数となるように。このことは最初に「gnomon」の理論に基づいており、これはユークリッド書IIにおける幾何学的な使用と一致する算術的意味を持っていた。形が gnomon のような図形を二つの正方形の側面に配置すると、それらはより大きな正方形になる。現在、数字1を一点で表すとします。この点の周りに3つの他の点を配置し、四つの点が正方形（1 + 3 = $2{2}$）を形成するようにします。その四つの点（正方形の隣接する二辺上）に等間隔で5つの点を配置すると、もう一つの正方形（1 + 3 + 5 = $3{2}$）が形成されます。そして続く奇数 1, 3, 5 ... を「gnomon」と呼び、一般式は

1 + 3 + 5 + ... + (2n - 1) = $n{2}$

である。次の奇数を加えると、すなわち 2n + 1 を加えると、$n{2}$ + (2n + 1) = (n +

1)2. すると、二つの平方数を求めるには、2n + 1が平方であることを確認すればよい。2n + 1 = $m{2}$とするとき、n = 1/2($m{2}$ - 1)であり、{1/2($m{2}$ - 1)}2 + $m{2}$ = {1/2($m{2}$ + 1)}2が得られ、mは任意の奇数である。これがピタゴラスに帰せられる一般公式である。

プロクルスもまた、ピタゴラスに帰す比例論理と五つの「宇宙的形体」（五つの正多面体）の構築法を記している。その中の一つ、十二面体は十二個の Pentagon で覆われており、正五角形の構築には「極端と平均比」（Eucl. II., 11, および VI., 30）が用いられる。これは Area の適用という方法の一例である。ピタゴラス自身にどれだけの部分が寄与したかは不明だが、この方法はピタゴラス派によって完全に開発され、幾何学的な有力な手法となった。

最も基本的なケースは Euclid, I., 44, 45において示されており、与えられた直線を底とする平行四辺形の作図法が述べられている。この構築法は算術的除算の幾何学的相当物である。一般化されたケースでは、平行四辺形は与えられた直線に適用され、その部分が与えられた直線を越えるか不足するかし、その余りまたは不足分が同じ角度と幅を持つ与えられた直線上の平行四辺形（底として）と相似である場合がある（Eucl. VI., 28, 29）。これは二次方程式 ax [+-] $mx{2}$ = C（実数解の場合）の最も一般的な形式の幾何学的相当物である。

古代ギリシャ幾何学の歩みにおいて、根が実数であることが解明された（ユークリッドVI，27）。重要なことに、この面積の適用法という手法は、パルガのアポロニウスによって、3つの円錐曲線の基本的な性質を導く際に直接用いられた。これらの性質は直交座標における円錐曲線の方程式と対応しており、アポロニウスが初めて与えたそれぞれの円錐曲線の名前は、理論から取られている。「パラボラ」（ギリシャ語：パラボレ）は「適用」を意味し（つまりこの場合、平行四辺形は直線に完全に適用される）、「ハイパロバ」（ギリシャ語：ハイポブレ）は「過剰」を意味し（つまりこの場合、平行四辺形は直線を超えて重なる）、「エッピロバ」（ギリシャ語：エルペプシス）は「欠如」を意味する（つまりこの場合、平行四辺形は直線に欠如している）。

また、ピタゴラス学派が解決した問題の一つは、与えられた直角多角形と相似な別の直角多角形と面積が等しい直角多角形を描くことである。プラUTERSキアは、この問題か、イウクリッドI，47の命題に基づいてピタゴラスが牛を犠牲にしたという説明があったことを述べている。

斜辺上の正方形の定理の主な特定的応用（例えばユークリッドIIの例）もピタゴラス学派によるものであり、与えられた長方形と等面積の直角多角形の構成（ユークリッドII，14）はその一つで、純粋二次方程式$x{2}$ = abの解法に対応している。

ピタゴラス学派は、任意の三角形の角度の和が二つの直角に等しいという定理を証明した（ユークリッドI，32）。

総じて言えることは、ピタゴラス幾何学はユークリッドの『要素』第一巻、第二巻、第四巻、第六巻の大部分をカバーしており、第六巻に関してはピタゴラス比例理論が共通量にのみ適用されたという限定がある。

古代ギリシャ幾何学の発展において、ユークリッド『要素』第3巻の命題に関する情報は不完全であるが、ヒポクラテス・チオス（紀元前470年頃から紀元前400年頃まで活動した）にその大部分が周知されていたことは確実であり、これより推測すると、第3巻の主要な命題もピタゴラスの幾何学には含まれていたことになる。

さらに、ピタゴラス派は無理数、すなわち不比長線を発見した。これはおそらく正方形の対角線が辺に対して不比長であることを初めて明らかにしたもので、その比率は√2:1であることが示された。この不比長性のピタゴラス派による証明はアリストテレスとユークリッド『要素』第10巻に挿入された命題によって保存されている。これは反証法により、対角線が辺と比で通じる場合、同じ数が奇数でありながら偶数であるという矛盾を示すものである。

この不比長の発見は幾何学者たちにとって大きな衝撃を与えた。なぜなら、これはピタゴラスによって作られた比率理論が普遍的に適用されないことを示し、それにより証明された命題が本当には確立されていないことを意味したからである。そのため、不比長の発見が秘密裏に保たれた時期があったという物語や、その最初の告発者が船難で死んだという話も存在する。

この幾何学的矛盾は、エウドックスス（紀元前408年から紀元前355年頃）によって大比例理論が発見され、『要素』第5巻で説明されているように、比で通じるものだけでなく不比長なものにも適用される理論により解消された。

ヒポクラテス・チオスの時代には、ギリシャ幾何学の範囲は『要素』に留まらず、直線や平面以外の特定の問題も取り組まれていた。

線と円についての古代ギリシャ幾何学の歩みは、立方体を二等分する問題、任意の角を三等分する問題、円を正方形に変換する問題など、これらの主要な問題を通じて最高の飛躍を遂げるギリシャ幾何学の方向性を決定することとなった。これらの問題に関する調査はヒッパルコスから始まり、《Elements》の完成と並行して進行した。

ヒッパルコス自身は、この二つの分野における同時研究の例である。一方では、彼はギリシャ人として知られる最初の人物で、 Elements の書籍を作成した人物であった。この本には、特にEuclidの『幾何原本』第3巻に含まれる円に関する重要な命題が含まれていたことだろう。しかし、もっと重要なものとして、ヒッパルコスは最初に円が直径の平方に等しい比で互いに比例することを証明したとされる。これは類似する円弧の部分も基底の平方に比例することから導かれる。これらの命題は、彼が『半月形の正方形化』という論文において使用し、その目的は円を正方形に変換することへの道を開くことだった。

後者の問題は、遠い昔から実用的な幾何学者にとって重要な課題であった。例えば、アナクサゴラス（紀元前500年頃-428年頃）は牢獄の中でこの問題に取り組んでいたとされる。ヒッパルコスの論文の主要な部分は、アリストテレスの『物理学』についてシムリキウスの記述において保存されており、その中にはエウデムスの『幾何学史』からの重要な断片が含まれている。ヒッパルコスは三つの異なる形態の特定の半月形を正方形化する方法を示し、最後に一定の円と一定の半月形の和を正方形化することも示した。不幸なことに、後者の半月形は正方形化できないものであり、その試みはそこで失敗したのである。

円を正方形に変形する試みは結局失敗した。

ヒポクラテスも立方体倍積の問題に取り組んだ。この有名な問題の起源には二つの説がある。一つには、古い悲劇詩人によってミノスが息子グラコスの墓標の大きさに不満を表明し、建築家に対してその大きさを2倍にすることを命じたという説がある。また別の説では、デリス人が疫病から救われるため、神託によりある立方体の祭壇を2倍にすることを求められたという。ヒポクラテスは確かにこの問題を解決しなかったが、二つの与えられた直線間の連続した比例における二つの中比を求める問題へと帰着させた。つまり、与えられた直線a, bに対してa : x = x : y = y : bとなるx, yを見つけることである。この条件はb/a = (x/a)3を満たすことが容易にわかる。特にb = 2aの場合には$x{3}$ = $2a{3}$となり、辺がaの2倍の立方体の辺長xを求めることができる。

その後、立方体倍積の問題は二つの中比の問題として唯一の形で試みられた。この問題に対する早期の重要な解決策は二つ記録されている。

(1) 公元前4世紀前半に活躍したタレンチュウム出身のアーキテュスは、三次元における非常に印象的な構成によって二つの中比を見つけ出した。この構成は実際には理論的であり、機械的な形式ではあったが、少なくともアーキテュスの手では立体幾何学は既に高度に発展していたことを示している。その構成は、空間内の一点を決定するために三つの表面の交点として定義された：（a）円柱、（b）円錐、（c）内部半径が0である「アンカーリング」。

(2) メネクムスは、エウドックスの弟子でありプラトンと同時代人で、双比例を円錐曲線を利用することで二つの方法で求めた。その方法は次の通りである。（α）二つのパラボラの交点を利用して、カートシアン座標系における方程式は \($x{2}$ = ay\) および \($y{2}$ = bx\) であり、（β）一方はパラボラと直角双曲線の交点を利用して、その方程式はそれぞれ \($x{2}$ = ay\) および \(xy = ab\) である。メネクムスがこの問題を解決する過程で円錐曲線を発見したことが推測される。これらの曲線はエラトステネスのエピグラムによれば「メネクムスの三重奏」であると呼ばれている。

角度の三分割は、イリスの sophist ヒッパシウスによって発明された曲線により達成された。ヒッパシウスはまたヒポクラテスやデモクリトスやソクラテス（紀元前470年から399年頃）と同時代人である。この曲線は「四分法」と呼ばれており、これは円を正方形に変換するための手段としても使用された（我々が知るところでは、メネクムスの兄弟ディノストラトゥスおよびニコメデスもこれを使用した）。理論的には、一定時間内に等速で動く二つの直線の交点を用いて構築される。具体的には、円の半径 OA と OB が直角に配置されているとする。A と B の接線が C で交わし、四角形 OACB を形成する。半径 OA は中心 O 回りに一定時間内に角度 AOB を描くように等速で動かされる。同時に AC は平行移動して等速で動き、A が線 AO 上を描く同じ時間内にその位置を移動する。これらの二つの動く直線の交点が様々な位置で形成される曲線が「四分法」である。

その後の幾何学に関する私たちの関心事は主に数人の人物によって行われたものであり、アブダラのデモクリトスとキレンのテオドロス（数学教師）などである。

プラト（前427-347）、テアテトゥス、エドックス、ユーメネウス、そしてユークリッドの時代に至る古代ギリシャ幾何学の歩みについて述べる。その名を聞く数学書の著者たちは以下の通りである。レオン（前408-355）はエドックスより少し年下で、ヒポクラテスによって収集されたよりも多くの有用な命題を集めた著作家であった。マグネシアのテウディウスはメナケムスとディノストラタスと同時期に生きたが、「部分的または限定的な命題をより一般的なものにまとめた」として、優れた数学書を作ったと推測される。テウディウスの本はアカデミーで使用された几何学教科書であり、アリストテレスも同様に使用した可能性が高い。

キレネのテオドロスとテアテトゥスは無理数論を一般化し、ユークリッドの第十巻（無理数に関するもの）の大半はテアテトゥスによるものであると考えられる。また、テアテトゥスは五つの正多面体（四面体、立方体、八面体、十二面体、二十面体）についても著作があり、ユークリッドの第十三巻（同様に五つの正多面体に関するもの）の内容もテアテトゥスによるものであると考えられる。第十二巻についてはエドックスが先駆者であり、プロクルスの言葉によれば、ユークリッドは要素を編纂する際にエドックスの多くの定理を集めた他、テアテトゥスによって完成された多くの定理をさらに完璧にし、前人による証明が不十分だった命題を不可欠な証明まで導き出した。

エドックス（約前408-355）はアーギメネデスの最も偉大な先駆者であり、特に「取り尽くし法」の発見に関心がある。天文学においては、エドックスは美しい同心円球体理論を考案して惑星の視運動、特にその視点と逆行を説明するために有名である。この理論は太陽や月についても適用される。

ユドックススはそれぞれの場合で、必要な球体を3つだけ使用しました。各惑星の運動は、4つの相互に連結した球体の回転の合成として表現されました。これらの球体はすべて地球の中心を通る直径を中心に回転します。最も外側の球体は日々の回転を表し、2番目の球体は黄道円または ecliptic 上の運動を示します；3番目の球体の極点は、黄道円上の2つの対立する点に固定され、その球体が2番目の球体の運動で回転します。そして、3番目の球体の表面には4番目の球体の極点が固定されています；4番目の球体は、その二極を結ぶ直径を中心に回転し、その球体に固定された惑星はその赤道上にある一点を通過します。極点、速度、および回転方向の選択により、惑星は実際には8字型の軌跡、つまり hippopede（馬縛り）と呼ばれる軌跡を描きます。この軌跡は黄道円に沿って配置され、その中心線で黄道円を二等分します。そして、その軌跡は黄道円と共に回転します。幾何学的な想像力の巧みな活用として、この仮説を類比するのは困難でしょう。

幾何学において、ユドックススは比例の偉大な理論を見つけ出しました。これは、共通比の量と非共通比の量にも適用可能な理論で、ユークリッドの『要素』第V巻で詳述されています。この理論は今でもその地位を保ち、永遠にその価値を持ち続けるでしょう。また、彼は特定の曲線を使用して2つの平均比例を解決しましたが、当時の文献にこれらの曲線の説明がないため、その性質については推測するしかありません。

最後に、最も重要なのは、彼が有名な取り尽くし法を用いて曲線の面積と立体の体積を測定することです。この方法の最も読者にとって馴染み深い例は、ユークリッド『要素』第12巻2章で証明された定理の証明でしょう。

円の面積は直径の平方に比例することである。この定理を含む証明は、すべての場合において、Euclidの書Xの命題1に相当する補題によって依存している。すなわち、二つの異なる種類の量があるとき、より大きいものからその半分未満を引くと、残りからはさらにその半分未満を引き続け、最終的に最初の小さい量よりも小さな量が残ることである。

アキメデスは、Euclid十二巻二命題の定理が、ある補題によって証明されたと述べている。すなわち、二つの異なる種類の量があるとき（すなわち、線分、面積、または立体）、大きい方から小さい方に比べて大きい差異があるとし、その差異は元の量と同じ種類の量を無限に加えることで任意の大きさを超えることができるという補題である。この仮定はアキメデスの公理や命題と呼ばれ、彼自身が述べているように、 exhaustion法を使用する人々によって以前から採用されていた。

実際には、Euclid十二巻二命題の証明に依存している Euclid十巻一命題の補題においても使用されており、Euclid第五巻第四定義と同様の声明であると考えられる。おそらくこの公理はEudoxusによって導入されたものであろう。

exhaustion法は一晩に発見されるものではなく、その方法を模索する痕跡を見ることができる。そのような方法が実際に形成される前に、アテネのソクラテス（紀元前470年-399年）と同時期にAntiphonという sophistが最初の一歩を踏み出した可能性がある。彼は円の中に正方形（または別の説では等辺三角形）を内接させ、その辺が描く弧を二等分し、多角形の辺数を倍増するように内接させる方法を使用した。彼はさらにこの過程を繰り返し、最終的に多角形の辺数を無限に増やすことで、円周に近づいていくと主張していた。

古代ギリシャ幾何学の歩みにおいて、多角形の辺を極めて小さくし、その多角形が円と一致するようにしたとき、この方法は形式的に誤りであったとしても、取り尽くし法の胚芽を含んでいた。

ヒポクラテスは、我々が見た通り、円の面積比が直径の平方比に等しいことを証明したと言われており、彼がこれを行ったのは取り尽くし法の何らかの形や予兆である可能性が高い。しかし、ユドックススによる役割については疑いの余地がない；彼は取り尽くし法を厳密な証明に基づいて基礎づけただけでなく、その方法を使って任意の台形体の体積（1）と円錐の体積（2）を求め、（1）任意の台形体は同じ底面を持つ同等の高さを持つ棱柱の三分の一であることを証明し、（2）任意の円錐は同じ底面を持つ同等の高さを持つ円柱の三分の一であることを証明した。しかしアリキメデスによれば、これらの二つの定理はデモクリトスによって初めて発見されたという驚くべき事実がある。彼はそれらを証明することができなかったが、おそらく完全な取り尽くし法による厳密な証明を与えることはできなかったことを意味する。

アリキメデスはこれらの定理に対するデモクリトスの貢献に大きな価値を認めるべきであると述べており、これは偉大なる人物が数学その他の分野においても驚異的な能力を持っていたという証拠でもある。アリストテレスの言葉によれば、「彼はすべてを考えたようなもの」である。

我々は他の情報からデモクリトスが無理数について書いたことを知っている；また、平行な円錐の断面（これらは明らかに非常に近接したものとされた）について議論したとも言われており、それらを不等または等しいと考えるべきか尋ねている。「もし」

古代ギリシャ幾何学の歩みにおいてアルキメデスに至るまで、不等であるならば円錐は多くの凹凸と段状の不均一性を持つ乱れた形となるが、もし等しいならば円錐は円柱の性質を有し、等しい円で構成されることが明白であり、これは非常に不合理である」という説明により、デモクリトスは無限小法に近づいていたことを示している。アルキメデスはさらに、直径の三倍比で球体が証明された定理は同一の補題を用いていると述べている。円錐、円柱、球体の体積に関する命題の証明は当然のことながらユークリッド『幾何学原論』第12巻に含まれており（それぞれProps.3-7 Cor., 10, 16-18）。

エウドックススの方法を例によって説明する必要がある。最も単純な例として、円錐に関する命題（ユークリッド『幾何学原論』第12巻、Props.10）を取り上げよう。与えられた四角形ABCDが円柱の底面であり、その底面と等しい円錐を持つとき、四角形ABCDを内接多角形として正方形とする。次に各辺によって生成される弧の中点を結び、8つの側を持つ規則的な内接多角形を描き、同様に16の側を持つ規則的な内接多角形を描き、そのような方法で繰り返す。各規則的な多角形に対して底面としてその多角形を基とする棱柱を立てることにより、円柱に内接する次々と増える高さが同じである棱柱を得る。每次棱柱の底辺の辺数を2倍に増やすごとに、円柱の体積より棱柱の体積から取り除かれるものは半分以上であり（円の面積と内接多角形の面積の差の半分以上を取り除くことによりEuclid XII., 2参照）、円錐の体積をV、円柱の体積をCとする。C = 3Vであることを証明しなければならない。もしCが3Vと等しくないならば、それは3Vより大きいか小さいかである。

まず（1）C > 3Vの場合を考える。そしてC = 3V + Eとなるように構造を続ける。

円柱内の棱柱を順に取り入れて、最終的な棱柱（体積 P）の外側に残された円柱の部分が全体の E よりも小さいとき、

C - P < E。

しかし C - 3V = E;

したがって P > 3V。

かつ以前の命題で証明済みであるように、P は同じ底面を持つ三倍の棱柱と等しい体積である。

従ってその棱柱は V（円錐の体積）より大きい：これは不可能であり、円錐は円錐体内包する棱柱を含んでいるからである。

したがって C は 3V よりも大きくない。

次に（2）C < 3V と仮定し、逆に V > 1/3 C。

この場合、円錐内の順次的な棱柱を取り入れて、円錐の外側に残された部分が V から 1/3 C の余分よりも小さいような棱柱まで到達するとき、

その棱柱は 1/3 C よりも大きいことがわかる。従って同じ底面を持つ棱柱（この棱柱は三倍の量である）を円柱内に取り入れたものは C よりも大きい：これは不可能であり、棱柱は円柱体内包されているからである。

したがって V は 1/3 C よりも大きくない、または C は 3V よりも小さくない。

従って C は 3V よりも大きすぎず、小さすぎず、等しいことになる；すなわち V = 1/3 C。

最後に追加するが、アルキメデスは主な放物線、双曲線、楕円の性質について十分に精通している。これらの性質は以前の論文で証明されており、アルキメデスはその論文を「放物線論」または「立体軌跡論」と呼んでいる。我々が知る二つの論文は（1）エウクリッドによる四巻の「放物線論」、（2）アリストイオスによる「立体軌跡論」であると推測される。これらの論文は失われており、前者は当然ながら、オッパロニウスの八巻の「放物線論」によって取り代わられた。

\clearpage

\chapter{第3章：アルキメデスの著作群と写本伝承}
\section*{THE WORKS OF ARCHIMEDES}

第三章。

アルキメデスの著作群。

アルキメデスの著書は、彼の Various treatises からその範囲を把握することができる。驚くほど多くの内容が彼自身の完全に新しい発見を反映している。彼は単なる編纂者や教科書作家ではなく、ユーメリウスやアポロニウスと異なり、彼らの著作は主に以前の幾何学者が使用した方法と得た結果を系統化し、一般化することに重点を置いていた。アルキメデスには既存の素材の単なる整理加工がない；彼の目標は常に何か新しいもの、知識総量への具体的な追加である。これは彼の各論文の前付された序文から確認できる。そこでは直接性、簡潔さと人間味が見られる。彼自身の発見と先駆者の成果との関係について、十分で公正な評価があり、自己中心的な影も見られない。彼は前任者によってなされた特定の発見が新しい方向へ拡張する可能性を提示したことを述べる。例えば、円を正方形に変換する試みにおいて先駆者の努力から、アルキメデスはパラボラシクセスを正方形に変換することを試みたが成功した。また、球体と円柱の体積と表面積についての彼の発見は、エウドックスによるピラミッド、円錐、円柱に関する定理の補足であると述べている。彼は特定の問題が長い間彼を難しめることも恐れずに述べており、ある場所では二つの命題について明確に述べ、更なる調査により誤りであることが判明したことを強調するために物語の教訓として挙げている。

アルキメデスの著作群と写本伝承

ギリシャ語版アーゲーミレスの普通の抄録では、彼の作品は以下の順序で記載されている。


1. 球と円柱論（二巻）。

2. 円周測定法。

3. 会計体と球状体論。

4. らせん論。

5. 平面均衡論（二巻）。

6. 沙算術。

7. パラボラの求積法。

最近では、このリストに重要な追加がなされた。1906年に、テキストの最新編集者であるジェンス・ルーカス・ハイベルガーは、コンスタンティノープルのパパドポロス・ケラメウスの jerusalemic 広場図書館で、数学的な内容を持つ pamphylos の写本を発見した。この写本は10世紀の良い手書きで Archimedes の著作が記載されていたことが判明した。古い文章を洗い出し（幸いにも不完全なものに留まった）、その皮紙は Euchologion として再利用された。しかし、多くの葉では、古い文章が多少読み取れる状態で残っている。この写本の重要な事実点は、以前失われたとされてきた二巻からなる論文 On Floating Bodies の大部分（1）と、メカニカル問題に関する Method という名前の書物の大半（2）を含むことである。幸運にも回復された重要な論文は、ハイベルガーの新しい（第二版）ギリシャ語版アーゲーミレスのテキストに含まれており（テューベナー、1910-15）、その概要は次の章で述べられる。

抄録に記載されている論文の順序は、作成時の順序ではなかった。しかし、 Various 前文と内部証拠から、著作の作成順序が推測できる。

アルキメデスの著作群と写本伝承について、次のような年代順序を確立できる：

1. 平面均衡論 I。
2. 抛物線の求積法。
3. 平面均衡論 II。
4. 方法。
5. 球と円柱論 I, II。
6. らせん論。
7. conoideとspheroides論。
8. 浮体論 I, II。
9. 円の測定法。
10. 沙塵計。

さらに、アラビア語経由で伝わってきた幾何学的命題の集めも存在する。これらは現在の形でアルキメデスによって書かれたものではなく、おそらく古代の著作を説明するために後のギリシャ人によって集められたものと考えられる。しかし、その中には特に優美な命題がいくつかあり、三半円と四半円からなる幾何学的図形に関するものなど、アルキメデスの起源にある可能性が高い。

また、アラビア語からの情報によれば、我々は三角形の面積を辺の長さで表す公式

\[
\Delta = \sqrt{s(s - a)(s - b)(s - c)}
\]

がヘロニウスによって幾何学的な証明が与えられたとされてきたが実際にはアルキメデスによるものであることを知る。

さらに、1773年にレッシングによって編纂されたギリシャ語のエピグラムによれば、アルキメデスはアレクサンドリアの数学者たちにエラトステネス宛ての手紙の中で有名な牛の問題を提示したと信用されている。プラトンの『 Charmides 』の注釈では、「アルキメデスの牛の問題」と呼ばれる問題についても言及がある。これは驚くほど難解なものである。

アルキメデスの著作群と写本伝承

アルキメデスの未発見された作品には、次のようなものがある。—

1. 多面体に関する調査はパッパスによって言及されており、パッパスは五つの正多面体について話した後で、アリクメデスによって発見された十三種類の半正多面体について説明している。これらの多面体は等辺と等角であるが相似ではない多角形で構成されている。少なくとも一つの半正多面体はプラトンによって既知であった。

2. 数学的内容を持つ『原理』という題名の書籍は、アルキメデス自身が学んだところによれば、数の命名について論じており、通常のギリシャ記数法では書き表せない大きな数を表現するシステムを説明している。『塵算術』（第五章参照）で同じシステムを展開した際、アルキメデスはそれを見たことがない人々のためにこれを説明することに利点があると述べている。

3. 『天秤に関する論文』（あるいは perhaps 『杠杆に関する論文』）。パッパスによれば、この著作では「同じ中心で回転するとき、大きな円が小さな円を制圧することが証明されている」とある。

4. 『重心に関する論文』はシムピリウスによって言及されている。しかし、この著作と前記の著作が別々の論文であることは確実ではない。第一巻『平面均衡論』はより大きな論文（おそらく『機械学要素』）の一部であった可能性があり、『天秤に関する論文』は別のタイトルであった可能性もある。『重心に関する論文』というタイトルは同じ論文を指すための緩い表現である可能性がある。

5. 『鏡像論』と呼ばれる光学的な著作があり、アレクサンドリアのテオンによって屈折についての言及が引用されている。

6. 『球体作成に関する論文』と呼ばれ、天体の運動を表すための球体の構築に関する機械工学的著作である（上記5-6参照）。

アラブの学者はさらにアルキメデスに属する作品を atributen おり、（1）円錐面についての論文。

アラビアの著述家たちは、さらにいくつかの著作をアルキメデスの作としている。すなわち、（1）『円について』、（2）『円に内接する七角形について』、（3）『互いに接する円について』、（4）『平行線について』、（5）『三角形について』、（6）『直角三角形の諸性質について』、（7）『所与命題集（Data）』の書である。しかしながら、これらの言及を裏付ける確たる証拠は現存していない。

\clearpage

\chapter{第4章：アルキメデスの幾何学 —— 取り尽くし法と求積法}
\section*{GEOMETRY IN ARCHIMEDES}

第四章。

アルキメデスの幾何学。

有名なフランスの数学者チャールズは、ユークリッドの二大後継者アリクメデスとアポロニウス・ペルガヌス（「偉大なる幾何学者」で、円錐曲線に関する古典的な論文を著した人物）の幾何学における主な特徴に、教唆的な区別を引き出した。チャールズは、これらの二人の作品が、幾何学において二つの大きな研究分野の起源と基礎であると述べている。アポロニウスは「形と状況の幾何学」に関心を持ち、一方でアリクメデスには「測量の幾何学」として、曲線平面図形の求積法と曲面の求積法および求積法と立方体積法についての調査があり、これらの調査はケプラー、カバリーニ、フェルマート、ライプニッツ、ニュートンによって考え出され、完成させられた無限に関する微分積分学を生み出した。

幾何学においてアリクメデスは、エウドックスの肩に立つかのように見えており、彼は新たなそしてより困難な求積法と立方体積法の新しいケースに対して Exhaustion 法を適用した。さらに、彼がこの方法を使用する際には、ユークリッド（おそらくエウドックスも同様）から知っているように、図形を「内接」するものだけではなく、「外接」する図形でも近似を行うという興味深い変化を導入した。ユークリッドと彼の推論法では、測定したい図形を「内接」する図形でしかなく、その二分の一の反証法を逆に使用して下からの近似が可能となった。一方でアリクメデスは、上からも下からも近似を行い、測定したい面積や体積に近づくために、内接する図形だけでなく外接する図形でもより密接な近似を行うことで、内接する図形と外接する図形を圧縮し、最終的にこれらが互いに一致し、測定したい面積や体積と一致することを目指す。

測量すべき図形。しかし、彼はギリシャ人が常に遵守した慎重な方法を採用し、特定の曲線や表面が内接または外接する図形の限界形であると明言することはありません；彼が主張するのは、我々は曲線や表面にどれほど近づくことができるかだけです。

取り尽くし法による論理的証明は、結果がどのように得られたかだけでなく、実際に行われた手順の真の性格さえも隠してしまう傾向があります。アリクメデスが実際に特定の場合で行うことは、解析的な等価物を記述したときに実質的な積分であることが明らかになる操作です。この言及は、放物線セグメントとらせんの面積、球体と球体セグメントの表面積と体積、二次曲線回転体の任意のセグメントの体積について適用されます。通常、これらの結果は、目的のために精密に鍛えられた一連の前提命題を経て得られます。この方法は、戦術家が予見し、直ちに有用でないものは排除し、順序に従って各位置を掌握し、最後に（スキームの詳細な展開が観客の心の中でその最終目的さえほとんど隠すまで）、突然一撃を与えるようなものと似ています。アリクメデスは、明確には直感できないと思われる命題を次々と提示しますが、後で必ず使用されることを知らされ、段階は容易に進むため、最初に提示された問題の困難さはほとんど認識されません。プラUTERSは言います、「幾何学においてより難しい問いや、より単純明瞭な問いを見つけることは不可能である」

アルキメデスの幾何学：取り尽くし法と求積法について論じる。しかし、プラUTERSクが次のように述べているように、その手法は明らかに言過りである。「我々は連続した手順の容易さによって、誰でもそれらを自発的に見つけ出すことができると錯覚する」。実際には、論文の洗練された単純さと完全な仕上げには同時に謎の要素が含まれている。各段階は前段階に依存しているにもかかわらず、アルキメデスがそれらをどのように提示されたのかについては明確ではない。ウォリスの言及にあるように実際には、「彼自身が調査の痕跡を隠蔽しようと努めたかのような印象を与える」という真実性がある。「彼は結果に同意を求めつつ、その手法の秘密を後世から奪うことに執着した」。

『方法』という著作では、部分的な例外が提供される。ここでは、アルキメデスの作業室の内部を見ることができるかのような veil を上げたような感覚がある。彼はどのようにして求積法と体積計算の定理を発見したのかを説明しており、同時に、（1）定理の真実性を提示するための手段が科学的な証明を提供しない場合でも、（2）承認される前に確認されるべき厳格な幾何学的手法による証明が必要であることを強調している。

アルキメデスはエラトステネスに宛てた手紙で述べている。「あなたが哲学において優れた学生であり、数学的探求を好む人間としてのあなたの地位を考慮に入れると、私はあなたに対してその特定の手法の特異性について詳細に説明する必要があると考えた。この方法を見れば、あなたは数学のいくつかの問題を機械的に調査する手段を得ることができるだろう。この手続きは、確実な幾何学的手法による証明が続く前に、その有用性が疑問視されないことを私は信じている」。

アルキメデスの幾何学：取り尽くし法と求積法

実際の定理の証明も含め、あることが機械的方法によって初めて明らかになった後で、そのことを幾何学的に証明することが必要となった。これは、その方法による調査が実質的な証明を提供しないからである。しかし、既に該当する問題について機械的方法により知識を得ている場合、証明を補足することは以前の知識なしで証明を見つけるよりも容易である。したがって、エウドックスが最初に発見した定理の証明、すなわち円錐は円柱の三分の一であり、棱柱と等しい底面を持つ正四角錐も同様であるという定理の場合、デモクリトスがこれらの図形についてこの真実を最初に主張したことには、大きな信用を与えるべきである。彼は証明しなかったが。

私は現在発表する定理の証明においても以前の発見と同様にその方法で発見した。そして、既に述べたことから虚言を被せられたくないとともに、数学にとって有用であると考えられるこの方法を記述し公開することにより、他の定理も新たに発見できるかもしれないという信念に基づいて、現在その方法を記述し公開することが望ましいと判断した。

まず最初に機械的方法によって初めて明らかになった定理を述べる。それは「直角円錐の断面の任意の部分は、同じ底面を持つ等しい高さを持つ三角形の四分之三である」という定理である。そして、その他の同様な方法で調査された定理もそれぞれ示す。

その後、

本邦書末において、Xの面積または体積を求めるための幾何学的証明を与える。

アキメデスが用いた機械的方法の大まかな特徴について、以下に説明すると希望される通り、概要を示す。Xが求められる面積または内容を持つ平面または立体図形とする。最も単純な場合においては、Xの無限小要素を既知の別の図形Bの同様の要素に対して秤で比較することにより、Xの面積または内容を求めることになる。Xの直径または軸を引くと、無限小要素は一般に軸または直径に直角に平行な断面であるが、必ずしもそうではない。したがって、すべての断面の重心は軸または直径上の一点または他の点にあり、その重さはそれぞれの点で作用することができる。平面図形の場合には無限小断面を「平行直線」と呼び、立体図形の場合には「平行平面」であるとし、無限数の断面の合計が全体のXを形成すると述べる。（これらの断面は実際には非常に狭い平面帯または非常に薄い板状体として捉えられるが、そのように表現される。）直径または軸は測定する図形から離れた方向に延長され、この直径または軸の延長を想像して天秤の棒またはレバーと見做す。現在の目標は、Xのそれぞれの要素をレバー上の一点で適用し、既知の図形Bの対応する要素は最初から実際に存在する点で作用することである。アキメデスは、したがって、測定する図形の各要素を移動させる方法を考え出すのである。

アルキメデスの幾何学において、Xの要素を元の位置からずらし、レバー上の一点に集中させることを考える。各要素が支点を中心に平衡するとき、Bの重心で作用する相当の要素についても同様である。XとBの要素は、支点からの距離に反比例する性質により、支点を中心に平衡している。B全体の重心は既知であり、その重心一処に一括りにして考えることができる。アルキメデスは、図形Bの各要素が実際に作用する点でそれぞれの「モーメント」の総和が、図形全体を重心一処に作用させたときのモーメントと等しいことを仮定している。また、Xのすべての要素は棒またはレバー上の固定点一処に集中される。この固定点をHとするとき、Bの重心をG、支点をCとするとき、X : B = CG : CHとなる。

したがって、Xの面積や内容を求めることができる。

逆に、面積や体積が既知の場合、この方法はXの重心を求めることにも使用できる。この場合、Xの要素とX自体は実際に作用する位置で作用し、既知の図形または図形の要素は固定点H一処で作用する必要がある。X、B、CHが既知であるため、比例式

B : X = CG : CH

から重心Gを求めることができる。

機械的方法は以下の（1）放物線部の面積、（2）球体と楕円体の体積、（3）球部の体積と直立部分のそれぞれの体積を求めるために使用される。

アルキメデスの幾何学：取り尽くし法と求積法

三次軸回転体、（4）球半分（a）の重心（a）、任意の球の部分（b）、任意の円錐体および放物面回転体（c）、半円柱（d）の重心（c）。すなわち、半円の重心。

アルキメデスは次に、この論文の特別な主題となる二つの立体の体積を求めることに進む。これらの立体は次の方法で生じる：--

（1）正方形底面を持つ直方体の中に内接する円柱があり、その円柱の両基底が対向する正方形面に内接する円であるとする。この円柱の一方の基底を含む正方形の一方の辺と、対向する基底の平行直径を通る平面を考える。この平面は、馬蹄のような形状を持つ立体を切り離す。アルキメデスは、このように切り離された立体の体積が直方体の体積の六分の一であることを証明する。

（2）正方形面に内接する円柱が立方体内に内接しているとする。また、対向する別の二つの面にも同様に関連した円柱がある。これらの円柱は、すべての角が丸みを帯びた立体を含んでいる。アルキメデスは、この立体の体積が立方体の体積の三分之二であることを証明する。

機械的方法でこれらの事実を証明した後、アルキメデスは取り尽くし法による厳格な幾何学的証明で論文を締めくくる。残念ながら、手稿は終盤に一部欠落しており、一定の修復が必要である。

次に、編纂版において最初に現れるアルキメデスの他の論文について簡潔に説明する。最初のものは次の通りである：--

『球と円柱に関する論文』。

書籍Iは、Dositheus（Cononの弟子）に宛てられた序文から始まります。その中で、彼は以前にDositheusに、直角円錐の「断面」（即ち、パラボラ）の任意の部分が同じ底と高さを持つ三角形の4分の3であることを証明した論文を伝達したことについて思い出させます。そして、彼はその後で発見したいくつかの定理の証明を送ることを追加し、これらの定理はエウドックスの台形体や円錐の体積に関する命題と比較しても価値があると考えていることを述べています。その定理は以下の通りです：（1）球面の表面積は四つの大きな円の面積に等しい（即ち、我々が「大円」と呼ぶもの）；（2）任意の球面部分の表面積は、頂点から底円上的一点までの直線で定義される円の半径に等しい円の面積に等しい；（3）球を内接する円柱があり、その高さが直径と等しい場合、（a）円柱の体積は球の1.5倍であり、（b）円柱の表面積（底面含む）は球の1.5倍である。

次にいくつかの定義が続きます。その後、特定の仮定が述べられ、その中には以下の2つが有名です：（1）同じ端点を持つすべての線の中で、直線が最も短い（これは直線の別の定義の基礎となっています）。（2）異なる線、面積、固体のうち大きいものが小さいものより、一定の量を繰り返し加えることで任意に大きな量を超えるような量を超越する（即ち、同じ種類の量）。これがアルキメデスの公理です。

彼はまた、線対（包括断線も含む）のペアについて仮定しています。

アルキメデスの幾何学：取り尽くし法と求積法

面積を測定する際に、同じ端点で限定され、同一方向に内接する凹曲面の対は、外側の方が内側よりも大きいという仮定がある。これらの仮定は彼の調査の基礎であり、その調査は常に内接図形と外接図形を使用して進行する。

前段階的な命題をいくつか述べた後、アルキメデスは円柱と円錐の表面積（1）を求めることに到達する（命題 13, 14）。その後、円柱と円錐に関するユークリッドの一部の命題を引用した後で、彼は本の主な業務へ移行し、球体と球冠の体積と表面積の測定を行う。

大きな円周に正多角形を内接させて回転させることにより、球体よりも大きい立体と小さい立体を得る。一連の命題で、内接図形と外接図形の（a）面積、（b）体積を求めることにする。次に、内接図形と外接図形が相似である場合、その図形の表面積は辺の二重比、体積は三重比であることを証明する（命題 32）。さらに、内接図形と外接図形の表面積と体積はそれぞれ球体の表面積と体積よりも小さいと大きいことを証明する（命題 25, 27, 30, 31 附則）。これで、取り尽くし法を適用して主な命題についての球体の表面積と体積に関する主な命題を証明できる材料が揃った。その後の本では同じ方法を使用する。

球の部分に対する手順は、半球未満の部分に対して回転曲面を内接させて外接させ、その部分の表面に関する定理は最終的に命題42で証明される。命題43では、半球より大きい部分の表面が導かれる。命題44では、任意の部分を含む球の扇形の体積が与えられる。

第2巻は、与えられた円錐や円柱と等しい体積を持つ球を見つけ出す問題から始まる；これは、二つの中間比の問題を解く必要があるため、その解決策が仮定される。命題2では、1と44を使用して、球の部分の体積に対する表現が導かれるし、命題3と4では、与えられた球を平面で切るとき（a）表面、（b）体積が与えられた比にあり得るようにする重要な問題が解決される。第二問題（命題4）の解法は困難である。アリクメデスはそれを直線ABを点Mで分割することにより解決し、

MB : （与えられた長さ）= （与えられた面積）: $AM{2}$

となるように分割する問題に帰着させる。

この問題の解法と可能性の限界は、アリクメデスによって発見され、ユトキオスによる注釈の中で保存されている断片によって与えられる。それはパラボラと直角双曲線の交点を使用して達成される。次に三つの構築問題が続き、最初の二つは与えられた球の部分と同じ相似比を持つ部分を構築することである。その体積や表面が別の与えられた球の部分と等しいようにする（a）体積、（b）表面。最後の二つの命題は興味深いものである。命題8では、平面で球を分割したときの二つの部分の体積V, V'および表面S, S'について証明され、より大きい部分に属するVとSが与えられる。

円の測量。

この論文は、私たちが受け継いだ形では、三つの命題のみを含んでいます。第二の命題は、第一と第三の命題から容易に導かれるものであり、その点で不適切な部分があります。

命題1においてアリクメデスは、正方形から始めて辺の数を順次倍加する方法で円内接多角形と外接多角形を描き、これらにより円の面積が直角三角形の面積に等しいことを従来の取り尽くし法によって証明します。命題3は有名な命題であり、アリクメデスは純粋な計算によって円周と直径との比（我々はこれをπと呼ぶ）の上界と下界を求めており、その結果は3-1/7 > π > 3-10/71となります。アリクメデスは正六角形から始めて辺の数を順次倍加し、96辺を持つ内接多角形と外接多角形まで到達します。円周の長さが二つの多角形の周囲長の間に存在することを見越して、アリクメデスは円の直径で表された二つの多角形の周囲長を計算します。彼の計算は、√3の値に対する近似（上界と下界）に基づいており、これは30度の余弦の逆数から始まります。アリクメデスは既知として265/153 < √3 < 1351/780を仮定しています。本文においては、計算手順の結果のみが記されています。

アルキメデスの幾何学：取り尽くし法と求積法には、いくつかの大数の平方根の近似値が与えられています。例えば、1373943-33/64の平方根は約1172-1/8、9082321の平方根は約3013-3/4、そして3380929の平方根は約1838-9/11とされています。このようにしてアルキメデスは、96辺の外接多角形の周長に対する円の直径の比として14688/(4673-1/2)を導き出し、これを3-1/7に丸めます。内接多角形に対応する値は6336/(2017-1/4)であり、アルキメデスはこれより3-10/71であると述べています。この例は、ギリシャ人がアルファベットによる数記法に恥ずかしがらなかったことを示しています。

\textit{円錐体と球面体論.}

献辞の中でアリストテレスがDositheusへ向けて述べている内容から、また内部の証拠からも推測できるように、この本全体はオリジナルであることが明らかです。アキメデスはまず、彼の円錐体と球面体とは何かを説明し、その後、それぞれの説明ごとに論文の主な結果を述べます。円錐体は二つあります。一つ目は「直角円錐体」で、これはパラボラ（「直角円錐の断面」という古い名称から派生した）という名前から導き出されたものであり、従って回転円パラボラ体です。二つ目は「鈍角円錐体」で、これは軸を軸として回転する双曲線（「鈍角円錐の断面」という名称から）によって生成された回転双曲円柱体です。球面体も二つあり、これらは楕円（「鋭角円錐の断面」という名称から）をその主軸と副軸についてそれぞれ回転させて生成される回転体で、一つ目は「長方」（または扁平）球面体、二つ目は「短方」（または延長）球面体と呼ばれます。これらの球面体の斜辺部分の体積は後に円錐体の体積を用いて求められます。

円錐体の断面を底とし、その頂点を頂点として、

そして、この底は斜交円錐の楕円截面であることにしたがって、アリキメデスは、底が楕円である円錐体を「円錐体の断片」と呼び、一般的な円錐とは区別する。

通常通り、いくつかの前駆的な命題が必要である。アリキメデスはまず、算術級数（a, 2a, 3a, ... na）とその項の平方（すなわち1², 2², 3², ... n²）を含む特定の級数の和を求め、幾何学的形式で表現する。これらの和は、各図形の無限個の要素に対する最終的な加算に必要なものであり、これは「積分」に相当する。

次に、双曲線と楕円に関する2つの性質（命題3）が述べられる。その後、斜交円錐の楕円截面による楕円の面積を取り尽くし法で決定する命題4が続く。さらに3つの命題があり、そのうちの最初の2つ（命題7, 8）は、先に述べた円錐体の断片は実際上斜交円錐の断片であることを示すものであり、これは実際に円截面を求めることによって行われる。命題9では、双曲面または球状体の楕円截面も特定の斜交円柱の截面であることを類似した証明がなされる（この円柱は、その軸が断片を切った双曲面または球状体の軸に並行である）。命題11-18では、各双曲面および球状体から断片を切り取る様々な截面の性質が示され、これらの截面は軸に対して垂直または斜交である場合によって円または楕円である。これらには接平面などのいくつかの性質も含まれている。

論文の本筋は、命題19, 20から始まる。ここでは、固体の断片から等間隔に平行な底面からの平面截面を多く引き出し、それらを描くことでどのようにして多くの要素の加算を行うかが示される。

円柱（一般に傾斜した）を各平面断面と交わす生成線が軸に対して平行であり、両側の隣接する断面で終焉するものにおいて、アルキメデスは、段状の階段のような外接図形と内接図形を作り出すことができる。これらの図形は、平行な面を持つ円柱の部分を構成し、円柱体の各部分の体積を求めることができる。これは取り尽くし法を使用して達成される（回転放物面の場合には命題21, 22、双曲面の場合には命題25, 26、球状体の場合には命題27-30）。結果は次のようである：（1）回転放物面上の任意の部分は、同じ底と軸を持つ円錐またはその部分の半分以上である；（2）回転双曲面上または球状体上の任意の部分は、同じ底と軸を持つ円錐またはその部分に対する比としてAD + 3CA : AD + 2CA（双曲面の場合）、および3CA - AD : 2CA - AD（球状体の場合）である。ここでCは中心、Aは部分の頂点であり、ADは部分の軸である（球状体の場合にはこの軸が半球状体より大きいとは限らない）。

らせんについて

アルキメデスに捧げられた序文は長く、その中ではまずコンオンへの記念を称え、次いで前二つの論文で含められた球と双曲面体および球状体に関する定理の要約が含まれている。その後、アルキメデスはらせんについて言及し、彼はこのらせんはこれまでのものとは異なる問題を提示し、共通点がないことを述べる。らせんの定義を与え、論文で証明される主な命題を述べた後、そのらせん（現在ではアルキメデスらせんと呼ばれる）は、ある点から等しい間隔で円周上に点が配置された点群の軌跡として定義される。

与えられた直線上の特定点（「原点」と呼ばれる）から始まり、その直線を一定速度で移動する点があり、同時にその直線は原点という固定点を中心に一定速度で回転している。定理1-11は前駆的なものであり、最後の2つは特定の級数の総和に相当し、これにより無限個の要素面積を加算する必要があるが、これは再び積分に相当し、曲線とその端点から引かれた2つの矢径で切り取られた図形の面積を求めることである。定理13-20は螺旋に対する接線の性質を興味深く難解な命題として確立している。定理21-23では、任意部分の螺旋に内接および外接する図形を示し、これらは原点を中心とする狭い円弧セクターからなる多様な図形であることを示す。これらの図形の面積は内接と外接する図形の面積の中間であり、通常の取り尽くし法によりアリッキメデスは必要な面積を求めている。定理24では螺旋の最初の一回り（r = aθの場合で1/3π(2πa)²）およびその部分（任意の点Pが最初の一回り上にある場合）の面積を与える。定理25, 26は同様に二回り目の螺旋と任意の弧（一回りを超えない）によって囲まれた面積について述べている。定理27では興味深い性質を証明しており、最初の一回りの螺旋で始線によって囲まれた面積を$R_1$、第二回りで追加されたリングの面積を$R_2$、第三回りで追加されたリングの面積を$R_3$とし、そのようなとき$R_3$ = 2$R_2$, $R_4$ = 3$R_2$, $R_5$ = 4$R_2$、そして$R_n$ = (n - 1)$R_2$であり、さらに$R_2$ = 6$R_1$であることを示している。

\textit{放物線の求積.}

この作品のタイトルは元々「放物線の切断」であったと思われる。

アルキメデスの幾何学：取り尽くし法と求積法

円錐において、直角を有する円錐はコニウムと呼ばれ、アポロニウスの時代以降に名称が変更された。この論文への前置きは、コニウスの死後、アポクリヌスとの最初の通信であることが明白であり、彼の亡き友人への悼辞的な言及から始まる。この前置きでは、アルキメデスが以前の幾何学者によって用いられた補題（アキメデスの公理、またはユークリッドX.1の定理）に基づく取り尽くし法について言及している。彼は、その補題を用いて証明された（1）円の面積が直径の二倍比に比例することと、球体の体積が直径の三倍比に比例すること、および（2）エウドックスによって証明された円錐と台形体の体積に関する命題について述べている。彼は、その直線と直角を有する円錐の断面で囲まれた部分（パラボラ）の正方形化がこれまで誰一人試みられていないことを述べる一方で、同一補題を用いてそのパラボラの部分が同じ底辺を持つ三角形の四分之三と等しいことを証明し、その証明は力学的調査から始めて幾何学的な証明まで行っている。この表現は、アポクリヌスへの前置きにおいても、メソッド同様にアルキメデスが機械的な調査を真実性の証拠と見なし、純粋な幾何学のみが絶対的な証明を提供することを示している。

パラボラに関する補題（一部は引用され、二つは明らかにオリジナルで証明される、命題4,5）についての力学的証明は17の命題まで続く。その他の補助的命題と共に幾何学的な証明が完成する。

アルキメデスの幾何学：取り尽くし法と求積法

（Props. 18-24）。機械的証明は、『方法』における証明とは一部類似しているが、より複雑である。パラボラの面積を測定する要素は直線ではなく細長い帯で構成されている。内接図形と外接図形はこれらの細長い帯からなり、鋸のような縁を持ち、すべての要素は台形であり、二つだけが三角形である。各台形（または三角形）は、ある固定点に対して吊られた他の面積と比較されて重さが測定される。したがって、内接図形と外接図形全体は、一地点から吊られた無数の面積の合計に重さが測定される。

結果は実際の積分によって得られ、通常通り取り尽くし法による証明で確認される。

幾何学的証明は次のように行われる。パラボラの部分において、同じ底と高さを持つ内接三角形を引いた後、残された各部分に同様の方法で三角形を内接させ、最初の内接三角形の1/4が新しい2つの三角形の合計であることを証明する。さらに、同じ方法で残された4つの部分に三角形を内接させ、それらの合計が前の2つの三角形の合計の1/4であり、従って最初の内接三角形の（1/4）2であることを証明する。このように進むと、パラボラの部分を尽きる一連の面積を得る。これらの合計は、最初の内接三角形を [Δ] で表す場合、

[Δ]{1 + 1/4 + (1/4)2 + (1/4)3 + . . .}

である。

アルキメデスは、Prop. 23において無限級数の総和が 4/3[Δ] であることを幾何学的に証明し、その後、この面積とパラボラの部分の面積が等しいことを反証法で確認する。

\clearpage

\chapter{第5章：砂の計算者（砂粒論） —— 巨大数と宇宙の大きさ}
\section*{THE SANDRECKONER}

第五章。

砂の計算者。

『砂の計算者』は独自の位置を占めています。数学的には重要な内容ではありませんが、アリ相信大家の多様性と天才を示す算術的奇蹟であり、ギリシャ天文学の歴史に関する貴重な詳細も含んでいます。これらの情報はアリシオン王ゲロンに捧げられた序文の中で述べられています。以下のように始まります：

「王よゲロン、ある者は砂の数が無限であると考えています。そして私は、その砂をシラクサやシカリアの他地域にも含めるとします。また、他の誰も住んでいない地域でも同様です。さらに、ある者は無限とは考えませんが、それでも何らかの数はその量を超えることはないと考える者もいます。そして彼らは、地球を中心とした宇宙という名前の球体を想像し、その中心点が地球の中心であるとします。この球体内に砂を含み、さらに海や地殻の空洞を砂で満たして最も高い山と同じ高さまで詰めると、それもまだ数え切れないほど多いと認識するでしょう。しかし、私はあなたに、幾何学的な証明を使って、私がゼウシッパスに送った著作の中で述べている数について示すことで、その砂の量を地球や宇宙の量を超えることが可能であることを説明します。

現在、あなたは「宇宙」が多くの天文学者によって最も高い山と同じ高さまで海と地殻の空洞を砂で満たした球体を指す名称であることを知っています。

半径は太陽の中心と地球の中心との直線に等しい。これは、あなたが天文学者から聞いた一般的な説明である。しかし、サモスのアリストテラコスは、前提に基づいて宇宙が現在呼ばれるものよりも遥かに大きい結論を導く仮説を含む本を出版した。彼の仮説は以下の通りである：固定星と太陽は動かない、地球は太陽を中心に円周上を回転し、太陽はその中心点に位置している、そして固定星の球面は太陽を中心に配置され、その大きさは地球が回る円の直径に対する比と同じであり、それが固定星までの距離に対する比と同じである。

したがって、アリストテラコスは、地球上の物質の量は宇宙の広大さに比べて微々たるものであるという絶対的かつ現代的な証拠を提供している。しかし、アキメデスの目的にはこれが適さないため、彼はアリストテラコスの仮説を一部修正する。彼は中心点（数学上の点）が球面の表面に対して何らの比も持たないと主張し、したがってアリストテラコスは、所謂「宇宙」の大きさに対する地球の大きさの比と同じように、新しい意味での実際の宇宙の大きさに対する所謂「宇宙」の大きさの比を意味していると解釈するべきである。

次に、彼は地球、月、太陽の一定の寸法を仮定し、地球の中心から太陽の直径が視覚的に形成する角度を推定しなければならない。そしてそれぞれの場合において、彼はその角度を過大評価する必要がある。

砂の計算者（砂粒論）：巨大数と宇宙の大きさ

確実な数字を用いて慎重に推定する。したがって彼は、地球の周囲が30万スタディアであることを証明しようと試みた人物もいるが、その数はエラトステネスが25万2千スタディアと推定し、約2万4千6百六十二マイルとすると直径は約7千8百五十マイルとなる。彼自身はそれを十倍に想定し、30万スタディアとする。

彼は続けて、地球の直径が月より大きいこと、太陽の直径が地球より大きいことを述べる。太陽の直径について、ユドックススはその直径を月の九倍と宣言したが、フィディアスは十二倍とし、アリストアクスは太陽の直径が月の十八倍以上二十倍未満であることを証明しようと試みた（この論文『太陽と月の大きさと距離に関する論説』は現在も存在しており、特定の仮定に基づいて厳密に結果を証明する優れた幾何学的作品である）。アリスタクスもまた慎重な推定を望むため、彼は太陽の直径を月の三十倍とし、それ以上にはしない。

さらに彼は、アリストアクスが発見したことを述べる。すなわち、太陽の直径が黄道円の約720分の1であることが見え、角度は約半度である。そして自己で作成した単純な装置により、太陽の直径が昇る直後における視点から見た角度を、約164分の1より小さく200分の1より大きいと求めた。観測者の地球表面からの視点での角度としてこの値を受け入れると、中心地心座標系において角度は約203分の1より大きくなることを幾何学的命題を用いて計算することができる。

したがって

太陽の直径は、所謂「宇宙」の大円に内接する正812角形の一辺より大きいことは言うまでもなく、さらに正千角形（1000辺の多角形）の一辺より大きい。

これらの仮定に基づき、正千角形（他の6以上の辺を持つ任意の正多角形についても同様）が大円に内接するとき、その周長は直径の3倍以上であることが容易にわかるため、地球の直径は1,000,000スタディオムより小さい一方で、所謂「宇宙」の直径は地球の直径の10,000倍未満であり、したがって10,000,000,000スタディオム未満である。

最後に、アキメデスは砂の量について仮定し、パプリコ種子の大きさよりも小さい砂の量は最多1万粒を含まないとし、パプリコ種子の直径は少なくとも40分の1のダクトリュス（スタディオムは10,000ダクトリュス未満）であると仮定する。

アキメデスは計算を行う準備ができており、しかしアルファベットによる数記法がこのような大きな数を表現するのに不十分であるため、彼は驚くべき大数表記の用語を開発する。

ギリシャ語には1万（10,000）までのすべての数に名前があり、従って通常の数記法で1万万（100,000,000）までのすべての数を表現することは難しくない。アキメデスは言いますように、これら全ての数を第一級の数と呼ぶことにしましょう。第二級の数は100,000,000から始まり、100,000,$000{2}$までで終わる。100,000,$000{2}$が第三級の最初の数であり、これは100,000,$000{3}$まで続く。そして、万万（100,000,000）の順序から始まり、100,000,000(100,000,000)まで続くもので、簡略化のためにPと呼びます。

すべての順序までPまでの数を第一期とし、第二期の最初の順序はPから100,000,000 Pまでとする。次に、第二期の第二順序はこれから始まり、第三順序は100,000,000² Pから始まり、同様にして第二期の100,000,000番目の順序まで進み、P²を含む1,000,000,000(100,000,000) Pで終わる。第一期の第三順序はP²から始まり、各順序は同じように進行する。一期ごとの系列とその各一期の順序を続けることで、最終的に第一期の100,000,000番目の期がP(100,000,000)で終わるまで到達する。この計画の驚くべき範囲は、最初の一期内の最後の数が800,000,000桁を持つ一方で、100,000,000期目の最後の数にはその8倍の桁が必要であることを考慮することで明らかになる。

アキメデスは「砂の数」というものを表現するために、最初の一期内の第八順序まで到達する必要はない。最初の一期内の各順序はそれぞれ1、$10{8}$、$10{16}$、$10{24}$、... ($10{8}$)(99,999,999)から始まり、すべての必要な数を10のべき乗で表現できる。

パピヨン種子の直径は少なくとも1/40ダクトゥス以上であり、球体はその直径の三倍比に比例する。したがって、直径が1ダクトゥスの球体は64,000個未満のパピヨン種子を含み、砂粒数も64,000 × 10,000以下であり、より一層言えば1,000,000,000、すなわち$10{9}$以下の砂粒数である。アキメデスは球体の直径を100倍に増加させ続け、それに応じた砂粒数を述べている。直径が10,000ダクトゥスの球体、より一層言えばスタディオンの球体には$10{21}$個未満の砂粒が含まれており、この方法で進めていくことにより...

直径100スタディア、10,000スタディアと続く円錐体において、彼は直径10,000,000,000スタディアの円錐体内の砂粒数に至り、これが所謂宇宙の大きさを表す。その対応する砂粒数は$10{51}$である。実際の宇宙の直径が所謂宇宙の10,000倍であるため、実際の宇宙内の砂粒数は$10{63}$となり、アキメデスの用語では第八位数の万万単位となる。

\clearpage

\chapter{第6章：力学 —— てこの原理・重心論と幾何学的証明}
\section*{MECHANICS}

第六章。

力学。

アキレストスは、数学的原理を用いて機械学を系統的に扱った最初の人物とされるが、彼によるそのような著作の痕跡は残っていない。実践的な力学においては、彼は飛ぶことができる機械鶴や子供たちを楽しませるための鳴る玩具を作成したと伝えられている（アリストテレスはこれを述べており、「それは子供たちが家の中での物損傷を防ぐために必要である」と加えている）。

アリストテレスの『機械学』においては、大きな重さを小さな力で動かす驚異的な現象についての言及があり、問題は様々な形の杠杆を通じてこの現象を説明する。また、力学における全ての動きがレバーとレバーマスコロス原理に帰着するとも述べられている（支点から力や重さが反比例関係にあるという仮定のもとで、それぞれの作用点からの距離が垂直方向に作用しているとされる）。しかし、レバーは「円周に参照されている」に過ぎない。より遠い位置で力が働く場合、それは大きな円を描くため、重さをより容易に動かすことができるという。

したがって、ここには証明がない。そのような性質を数学的に証明したのはアキメデスである。彼は特定の公理に基づいて、学習者の信頼を大きく要求しないで、レバーマスコロス原理を証明したのである。『平面均衡論』というタイトル通り、この著作は静力学のみに焦点を当てており、第1巻の命題6,7でレバーマスコロス原理が確立された後、残りの部分では特定の図形の重心を求めることに専念している。この著作には動力学は含まれていないため、速度平行四辺形も考慮されていない。

アリストテレスの『力学』において比較的適切な証明がなされている。

アキメデスの公理には、次のような仮定が含まれている：（1）等しい重さが等しい距離に置かれた場合、それは均衡し、等しい重さが異なる距離に置かれた場合、それは均衡しないが、「大きな距離にある重さの方へ傾く」という言葉で表現され、つまり、大きな距離にある重さの作用はその方向に運動を生じさせる；（2）および（3）重さが均衡している状態において何かが一方の重さに加えられたり取り除かれた場合、それは「傾く」方向へ向かい、加えられた重さの方へ傾くか、取り除かれた重さから何も取り除かれた方へ傾く；（4）および（5）等しい相似な図形が一致する部分を合わせて配置された場合、その重心も一致し、不等しい相似な図形が配置された場合、その重心は図形に対して同一の位置関係にある。

主たる命題である、二つの量が反比例する距離で均衡することについて、まず有理的な量から無理な量へと証明される。前段階の命題では、等しい量が直線上に等しい距離に配置され、奇数または偶数個の場合について述べられており、全体の重心がどこにあるかを示している。まず有理的な量の場合を考える。A, B をそれぞれ E, D に作用する重さとし、ED の C 点で A : B = DC : CE となるように ED を分ける。アキメデスは、このシステムが C において均衡することを証明しなければならない。ED を K まで延長して DK = EC、さらに DE を L まで延長して EL = CD とする。すると LK は C を中点とする直線となる。また、LK 上の H 点を LH = 2LE または 2CD のように取ると、残りの HK = 2DK であることがわかる。

トーマス・リトル・ヒース著『アキメデス』より力学の部分、特にてこの原理と重心論に関する幾何学的証明について記す。

2EC. A, Bが共通比であるとき、EC, CDも共通比である。xをEC, CDの共通比とする。wをAの一部としてxと同じ比率を持つ重さとする。すると、wはBの一部としてxと同じ比率を持つことになる。アキメデスはLH, HKをxで等分し、A, Bもwで等分する。さらに各wをそれぞれのxの中央点に配置する。これによりすべてのwはCを中心に均衡している。しかし、LH上の各点でのwの作用は全体としてEにおいて働くAと同等であり、同様にHK上の各点でのwの作用はDにおいて働くBと同等である。したがって、E, DにおけるA, BはCを中心に平衡している。

7節では、A, Bが共通比でない場合でも同じ結果を反証法により導く。8節では、全体と部分の重心が既知の場合、余分な量の重心を求めることについて示す。9-15節では、平行四辺形、三角形、平行台形の重心を求めることについてそれぞれ説明する。

第2巻は十節で構成され、主に放物線部分の重心を求めることに専念している。これは優雅だが難しい幾何学的作業であり、通常取り尽くし法により確認される。

\clearpage

\chapter{第7章：流体静力学 —— 浮体の原理とエウレカの発見}
\section*{HYDROSTATICS}

第７章。

流体静力学。

流体静力学は、静力学よりもさらにアキメデスの原創的な創造である。アキメデスの著作『浮体論』においては、彼に先駆者がないように思われる。彼の作品の中で証明された唯一の事実は、アリストテレスによって一種の証明が与えられているが、それは静止した流体の表面が地球の中心を中心とする球面であるという命題である。

アキメデスは、その理論全体を2つの公理に基礎づけている。これらの中の1つは最初から始まり、もう1つは第1巻の命題7の後に述べられている。公理1は次の通りである：

「流体が、その部分が均等に配置され、連続している場合、より圧縮された部分は圧縮されていない部分を押し出し、各部分は上部から垂直方向に圧縮されるという性質を持つと仮定する。ただし、流体が何かによって閉じ込められて圧縮されている場合はこの限りではない。」

公理2は次の通りである：「任意の物体が水の中で浮かぶ場合、その重心を通る垂線を通過して物体は浮かび上がる。」

命題2では、アキメデスはどんな流体でも静止した場合の表面が地球の中心を中心とする球面であることを証明している。

命題3から7まででは、固体（1）流体と同じ重さ、（2）流体よりも軽い、（3）流体よりも-heavyな状態での挙動について述べている。流体よりも軽い固体は完全に沈められず、その代わりに浮かび上がり、沈められた場合、上昇力は固体の重さと流体が排水された部分の重さとの差に等しいということが証明されている（命題5、6）。流体よりも-heavyな固体は、沈み続ける。

流体静力学において、固体が液体中に置かれた場合、その固体は底まで沈み、その固体の重さは液体に浸したときの重さより軽くなる。これは、液体を排水する量分だけ軽くなる（命題7）。

この最後の定理は自然と、ヒロニウムのために作られた冠に関する物語に関連付けられる。その冠が完全な金ではなく、銀も混ざっている可能性があると疑われ、ヒロニウムはアキメデスに金属の比率を決定する問題を提示した。この問題の解決策を見つけるとき、アキメデスは風呂場で裸で走り去り、[ギリシャ語: エウレカ、エウレカ]と叫んだという話がある。定理の一つを使用して解を導いたという説明もあるが、全体として見れば、実際の発見はビトゥーヴィウスによって記述されたより基本的な方法でなされた可能性が高い。彼は、風呂場に水が満ちているときに足を踏み入れると、その体積分量の水が溢れ出ることを見たとされるが、この考えを冠の場合にも適用し、（1）冠自体が排水する液体の体積V、（2）純金の同等重さが排水する液体の体積V1、（3）純銀の同等重さが排水する液体の体積V2を測定することに思い至った。

したがって、V = w1/W . V1 + w2/W . V2 となる。従って、WV = w1V1 + w2V2。

流体静力学において、冠の中に含まれる金と銀の重量比を求めることを考える。すなわち、

（w1 + w2）V = w1V1 + w2V2,

よって、

w1/w2 = (V2 - V)/(V - V1),

これにより、冠の中に含まれる金と銀の重量比が求められる。

第I書の最後二つの命題では、球体の断片が流体中に浮かぶ場合を調査する。その底面が（1）完全に表面以上であり（2）完全に表面以下である場合について考察し、どちらの場合も軸が垂直である位置で平衡していることが示される。この平衡は最初の場合において安定である。

第II書は純粋な幾何学的な傑作である。アキメデスは、回転放物面の直角断片が底面を上向きまたは下向きに浮かぶ場合（完全に表面以上または完全に表面以下）の数々の状況について調査する。その位置が平衡であり安定である場合において、固体の軸が垂直でないとき、その軸が垂直に対して傾斜する角度は完全に決定される。

比重の概念は通じているが、この実際的な用語は使用されていない。アキメデスは固体が流体よりも軽い、または重い、または等しい重量であることを述べる。また、重量比を表す必要があるときには、等量の体積に対する固体の重量に対して流体の重量との一定の比率を持つ固体について述べている。

\clearpage

\chapter*{付録A：参考文献解題}
\addcontentsline{toc}{chapter}{付録A：参考文献解題}

参考文献解題。

アキメデスの著作とユトシウスの注釈を含む最初の版は、1544年にバーゼルでヘラヴァージュ（ヘルワゲン）によって刊行された。D. リヴォア（パリ、1615年）はギリシャ語での命題と改訂されたラテン語証明を提供した。砂積算 ($Arenarius$) および円周の測定 ($Dimensio circuli$) とユトシウスの注釈は、ウォルズィが1678年（オックスフォード）でラテン語訳と注釈を付けて編纂した。トーレリの巨著（オックスフォード、1792年）では、ギリシャ語版の全著作とユトシウスの注釈、新しいラテン語訳が含まれており、最近まで標準的な文書として使用された；現在ではヒルベルグによって編纂され、3巻（テューベナー、レipzig、初版1880-1年；2版、新たに発見された $Method$ 等を含む1910-15年）で定義的な文書がその地位を取っている。

翻訳については、次のようなものがある。1558年に刊行されたアルディン版（4度目）には、コンマンドィヌスによる \textit{円周の測定} 、\textit{らせん} 、\textit{パラボラの求積法} 、\textit{軸体と球体に関する論文} 、$砂積算$ のラテン語訳が含まれている。アイザック・バローの版は、1675年に刊行された $アキメデス著作集$ 、\textit{アポロニウスペルゴイの円錐曲線論文集} 、\textit{テオドシウスの球面幾何学} 、\textit{新しい方法で解釈され証明される} に含まれている。作品の最初のフランス語版は1808年にピヤールによって2巻で刊行された。価値のあるドイツ語版（ノート付き）は1824年にストラルズブルグでエーネスティン・ニッツェによって刊行された。現代記号を使用した完全な版は、T. L. ヘイチ ($アキメデス著作集$ 、カムブリッジ、1897年）と \textit{アキメデスの方法} （カムブリッジ、1912年）で補完されたものがある。

\clearpage

\chapter*{付録B：アルキメデスと古代科学史年表}
\addcontentsline{toc}{chapter}{付録B：アルキメデスと古代科学史年表}

年代順序。

一部は正確でないものも含まれている。

紀元前

624-547  テラス

572-497  ピタゴラス

500-428  アナガゴラ

470-400  / キシオスのヒッパルコス

\ エリスのヒッパラクス

460-385  ケレネーのテオドロス

430-360  テラトゥムのアーキユタス

427-347  プラトン

415-369  テアエテトゥス

408-355  クニドスのユドックス

/ レオン

活動年代は約350年頃 < メネアクム

| ディノストラタス

\ テュディウス

活動年代は300年頃    ウォリドゥス

310-230    サモスのアリストアクス

287-212    アーキメデス

284-203    エラトステネス

265-190    ペルガマのアポロニウス

\clearpage

\end{document}
